% ============================================================
% Page 7 — Alban aux « siècles bénédictins »
% Transcription du document source.
% Texte conservé sans reformulation ni correction.
% ============================================================

\thispagestyle{plain}

\vspace*{12mm}

\begin{center}
    {\Large\bfseries Alban aux « siècles bénédictins »}
\end{center}

\vspace{18mm}

\noindent
Le 19 juillet 1228, entre dans sa « ville d'Albi » son nouvel évêque,
Durant de Beaucaire que Simon de Sully, archevêque de Bourges et
primat d'Aquitaine, appelle « homme de culture, de bonté, d'honneur ».
C'est grâce à lui qu'Alban entre dans l'histoire avec un document écrit
conservé aux archives départementales du Tarn.

\par
En effet, un de ses premiers actes fut de réorganiser le diocèse et
de faire entrer sous son obédience de multiples paroisses réparties
en huit archiprêtrés, dont celui de Négremont. Par accord passé en
1231 avec Adhémar, Abbé de St Benoît de Castres, nous apprenons
qu'une partie de la paroisse de Nougaret, et que la totalité de celle
de Sainte Marie d'Alban, demeurent sous juridiction bénédictine.
C'est donc l'abbaye qui pourvoit au prieuré de Sainte Marie d'Alban.
L'évêque d'Albi n'a d'autre droit que de se faire héberger pour la
visite pastorale ou de percevoir une taxe équivalente appelée
« synodaticum ». « Sancta Maria de Albanio » reste administrativement
compris dans l'archiprêtré de Négremont. Le dernier archiprêtre de
Négremont sera Guillaume de Pressoris. Le pape Eugène IV confirme
le 8 janvier 1438 la décision de l'évêque d'Albi, Robert Dauphin,
supprimant le susdit archiprêtré. Les revenus de la susdite paroisse
serviront à l'entretien de la maîtrise de Sainte Cécile d'Albi.
Alban deviendra chef lieu de claverie, ... mais le douzième siècle
va faire auparavant parler de lui avec l'hérésie cathare. Alerte aux
Albigeois !

\vfill
